%======================================================================
\section{Concluding remarks}
\label{sec:concl}
%======================================================================

We have constructed the soft light-cone wave function and the corresponding unitary evolution
operator $\OM$ of a fast-moving hadron explicitly through $\order{g^{2}}$ in QCD with $\Nf$ massless
quark flavours, within the CGC framework, extending the pure Yang--Mills construction of Paper~I.
Starting from the light-cone QCD Hamiltonian in the eikonal and Born--Oppenheimer approximations,
we showed that the valence sector couples to the soft modes only through its total colour charge
density $\rho=\rho_{g}+\rho_{q}$, which obeys the same non-Abelian algebra as in the pure gauge theory,
while the soft quark degrees of freedom generate new operator structures in $\OM$. The evolution
operator was organized as a fully normal ordered series in the soft gluon, quark and antiquark
creation and annihilation operators, with coefficients that are non-commuting functionals of $\rho$.

The coefficients were determined by the same two complementary procedures as in Paper~I. The
unitarity condition $\OM^{\dg}\OM=1$, imposed order by order in the coupling, provided the relations
among the coefficients, and the coefficients of real transitions were fixed by matching the action of
$\OM$ on the vacuum, one, two and three-gluon states, a single quark and antiquark, a quark--antiquark
pair and a quark--gluon state to the LCWFs computed in LCPT. Three features of the construction are
specific to the quark sector. First, a complete ansatz must contain, besides the structures that
create and annihilate soft pairs, the coefficients $\mf F$ and $\mf G$ for gluon emission by soft quarks
and antiquarks, which are of order $g$, and the coefficients $\mf P$, $\bar{\mf P}$ for the eikonal
rescattering of a soft quark on the valence charge, which are of order $g^{2}$; without them $\OM$ is
neither unitary nor able to diagonalize the Hamiltonian. They enter the unitarity relations of the
pair coefficients through the $t$- and $u$-channel quark and gluon exchanges. Second, in the
coefficients $\mf A_{1}$ and $\mf P$ the instantaneous interaction with the valence charge cancels
exactly the helicity conserving part of the corresponding one gluon exchange, as required by
unitarity. Third, every second order coefficient is a sum of products of first order coefficients
weighted by ratios of energy denominators plus instantaneous terms, Eq.~\eqref{eq:reduction}; for
independent subprocesses the overall energy denominator cancels and the coefficients factorize
exactly, while on-shell poles survive only in the particle number preserving coefficients and in the
$2\to2$ scattering coefficients, where they are defined by the M{\o}ller $i\epsilon$ prescription.

As an application we used $\OM$ to diagonalize the soft QCD Hamiltonian. At order $g$ the
Hamiltonian is diagonalized exactly to its free part once the coherent operator and the three-gluon
coefficient of Paper~I are supplemented by $\mf A_{3}$, $\mf F$ and $\mf G$. At order $g^{2}$ all quark
induced off-diagonal terms are removed; the quark loop contribution to the gluon self-energy and the
quark self-energy are scaleless and vanish in dimensional regularization, the ultraviolet pole of the
gluon normalization being proportional to $11N_{c}-4T_{F}\Nf$; and the only diagonal correction is the
coherent background field energy $-\int_{p}g^{2}\rho^{a}(p)\rho^{a}(-p)/(p^{+})^{2}$, now with $\rho$ including
the valence quark charge.

The construction presented here provides the full QCD evolution operator needed for the computation
of physical observables at next-to-leading order with quarks. In particular, together with the
order $g^{3}$ coefficients it is the input for the single inclusive gluon spectrum at NLO at
mid-rapidity in full QCD \cite{LRS:inprep}, for the inclusive production of soft quarks and
quark--antiquark pairs \cite{Blaizot:2004wv,Kovchegov:2006qn}, and for the quark contributions to
multi-particle correlations \cite{Kovner:2011pe,Altinoluk:2018byz}. We leave these applications for
future work.

\acknowledgments

The author is grateful to Michael Lublinsky and Vladimir Skokov for constant encouragement to
continue this work and to Tiyasa Kar for cross checking critical parts of the calculation. The author
also thanks the Department of Physics at North Carolina State University for providing research
assistantship support during the course of this work. This work is supported by the U.S. Department of
Energy, Office of Science, Office of Nuclear Physics through Contract No. DE-SC0020081.
